\section{System Design and Methodology}
\label{sec:method}

\subsection{Design Goals}
\cs\ was built around four goals. \textbf{(G1) No network egress:} a
complete audit must run with the machine's network interface disabled, so
that no code, prompt, or result is ever transmitted. \textbf{(G2)
Multi-modal:} pattern analysis, structural metrics, semantic (LLM)
review, and dynamic execution should contribute to one result rather than
living in separate tools. \textbf{(G3) Context-preserving triage:} every
finding must be inspectable next to the exact source that produced it.
\textbf{(G4) One number, honestly derived:} the platform should output a
single interpretable health score whose computation is fully disclosed.

\subsection{Architecture}
\label{sec:arch}
\begin{figure*}[t]
  \centering
  \begin{tikzpicture}[
      font=\footnotesize,
      chip/.style={rounded corners=2.5pt, draw=csslate!45, fill=white,
                   line width=0.5pt, minimum height=7mm, minimum width=30mm,
                   align=center, font=\scriptsize, inner sep=3pt},
      lyr/.style={font=\small\bfseries, color=csink, anchor=east,
                  align=right, text width=24mm},
      dn/.style={-{Stealth[length=2mm]}, draw=csslate!60, line width=0.8pt},
    ]

    % ---------- Presentation ----------
    \node[chip] (p1) {React renderer};
    \node[chip, right=4mm of p1] (p2) {12 coordinated screens};
    \node[chip, right=4mm of p2] (p3) {offline status strip};
    \node[lyr] at ([xshift=-6mm]p1.west) {Presentation};

    % ---------- Orchestration ----------
    \node[chip, below=10mm of p1] (o1) {Electron main};
    \node[chip, right=4mm of o1] (o2) {typed IPC bridge};
    \node[chip, right=4mm of o2] (o3) {offline policy};
    \node[lyr] at ([xshift=-6mm]o1.west) {Orchestration};

    % ---------- Services ----------
    \node[chip, below=10mm of o1, minimum width=21mm] (s1) {RepoService};
    \node[chip, right=3mm of s1, minimum width=21mm] (s2) {AiService};
    \node[chip, right=3mm of s2, minimum width=27mm] (s3) {ExecutionService};
    \node[chip, right=3mm of s3, minimum width=21mm] (s4) {RiskModel};
    \node[lyr] at ([xshift=-6mm]s1.west) {Services};

    % ---------- Local infrastructure ----------
    \node[chip, below=10mm of s1] (i1) {SQLite store};
    \node[chip, right=4mm of i1] (i2) {Docker engine};
    \node[chip, right=4mm of i2, fill=csgreenbg, draw=csgreen!45] (i3)
      {Ollama container \textbf{Llama 3.2}};
    \node[lyr] at ([xshift=-6mm]i1.west) {Local infra.};

    % ---------- layer backdrops ----------
    \begin{scope}[on background layer]
      \node[rounded corners=4pt, fill=csbluebg,  fit=(p1)(p3), inner sep=3mm] {};
      \node[rounded corners=4pt, fill=csgraybg,  fit=(o1)(o3), inner sep=3mm] {};
      \node[rounded corners=4pt, fill=csgraybg,  fit=(s1)(s4), inner sep=3mm] {};
      \node[rounded corners=4pt, fill=csgreenbg, fit=(i1)(i3), inner sep=3mm] {};
    \end{scope}

    % ---------- flow ----------
    \draw[dn] (p2.south) -- ++(0,-3mm);
    \draw[dn] (o2.south) -- ++(0,-3mm);
    \draw[dn] (s2.south) -- ++(0,-3mm);
    \draw[{Stealth[length=2mm]}-{Stealth[length=2mm]}, draw=csslate!50,
          line width=0.7pt]
      ([xshift=-30mm,yshift=4mm]p1.west) -- ([xshift=-30mm,yshift=-4mm]i1.west)
      node[midway, sloped, above, font=\scriptsize\itshape, color=csslate]
      {request / result};

    % ---------- no-egress boundary ----------
    \begin{scope}[on background layer]
      \node[draw=csslate!70, dashed, line width=0.9pt, rounded corners=6pt,
            fit=(p1)(p3)(i1)(i3), inner sep=8mm] (box) {};
    \end{scope}
    \node[anchor=west, font=\scriptsize\bfseries, color=csslate,
          fill=white, inner sep=2pt]
      at (box.north west) {\,Local machine\ \ $\bullet$\ \ no network egress\,};
    \node[anchor=east, font=\scriptsize, color=csred, fill=white, inner sep=2pt]
      at (box.north east) {\,$\varnothing$ network\,};

  \end{tikzpicture}
  \caption{\cs\ architecture. Four in-machine layers inside a hard
           no-egress boundary: a React presentation layer, an Electron
           orchestration layer, four analysis services, and local
           infrastructure hosting SQLite, the Docker engine, and the
           quantised \mbox{Llama~3.2} model. No stage opens an outbound
           connection.}
  \label{fig:architecture}
\end{figure*}


\cs\ is an Electron~\cite{electron} desktop application with a strict
main/renderer split. \Cref{fig:architecture} shows the four layers, all
resident on the developer's machine inside a no-egress boundary.

\begin{itemize}
 \item The \emph{presentation layer} is a React single-page application
 of twelve coordinated screens (\cref{sec:interface}). It holds no
 analysis logic; it issues typed requests over an
 inter-process-communication (IPC) bridge and renders results.
 \item The \emph{orchestration layer} (the Electron main process)
 validates each request, enforces the offline policy, and routes
 work to the services.
 \item The \emph{service layer} contains four in-process modules: a
 \textsc{RepoService} (checkout, file inventory, pattern
 detection, complexity), an \textsc{AiService} (prompt assembly,
 model I/O, finding extraction), an \textsc{ExecutionService}
 (Dockerfile synthesis, image build, sandboxed run, telemetry),
 and a \textsc{RiskModel} (fusion and scoring).
 \item The \emph{local infrastructure layer} comprises an embedded
 SQLite database for projects and findings, the host Docker
 engine, and an Ollama container hosting the quantised
 \mbox{Llama~3.2} weights~\cite{ollama,touvron2023llama2}.
\end{itemize}

\subsection{Analysis Pipeline}
\label{sec:pipeline}
\begin{figure*}[t]
  \centering
  \begin{tikzpicture}[
      font=\footnotesize,
      node distance=6mm,
      stage/.style={rounded corners=3pt, draw=csslate!45, fill=csgraybg,
                    line width=0.6pt, minimum height=9mm, minimum width=20mm,
                    align=center, inner sep=4pt},
      ana/.style={rounded corners=3pt, draw=csblue!45, fill=csbluebg,
                  line width=0.6pt, minimum height=8mm, minimum width=30mm,
                  align=center, inner sep=3pt, font=\scriptsize},
      out/.style={rounded corners=3pt, draw=csgreen!50, fill=csgreenbg,
                  line width=0.6pt, minimum height=9mm, minimum width=20mm,
                  align=center, inner sep=4pt},
      fl/.style={-{Stealth[length=2mm]}, draw=csslate!60, line width=0.8pt},
    ]

    \node[stage] (repo)  {Repository};
    \node[stage, right=of repo] (inv) {File\\inventory};

    % parallel analyses
    \node[ana, right=14mm of inv, yshift=10mm]  (pat) {Pattern scan\\\scriptsize(rules, CWE)};
    \node[ana, right=14mm of inv]               (cx)  {Complexity engine\\\scriptsize(McCabe $CC$)};
    \node[ana, right=14mm of inv, yshift=-10mm] (llm) {LLM review\\\scriptsize(Llama 3.2, schema)};

    \node[stage, right=14mm of cx] (fuse) {Finding\\fusion};
    \node[stage, right=of fuse]    (dyn)  {Dynamic\\sandbox};
    \node[stage, right=of dyn]     (risk) {Risk\\model};
    \node[out,   right=of risk]    (rep)  {Report\\\scriptsize PDF / JSON};

    \draw[fl] (repo) -- (inv);
    \draw[fl] (inv.east) -- (pat.west);
    \draw[fl] (inv.east) -- (cx.west);
    \draw[fl] (inv.east) -- (llm.west);
    \draw[fl] (pat.east) -- (fuse.west);
    \draw[fl] (cx.east)  -- (fuse.west);
    \draw[fl] (llm.east) -- (fuse.west);
    \draw[fl] (fuse) -- (dyn);
    \draw[fl] (dyn)  -- (risk);
    \draw[fl] (risk) -- (rep);

    \node[font=\scriptsize\itshape, color=csslate, above=1mm of pat] {in process};
    \node[font=\scriptsize\itshape, color=csslate, below=1mm of llm]
      {local container};

  \end{tikzpicture}
  \caption{The \cs\ analysis pipeline. A file inventory feeds three
           parallel analyses---pattern scanning and complexity in
           process, semantic review by the on-device model---whose
           output is fused (\cref{sec:fusion}), augmented by a sandboxed
           dynamic pass, and aggregated into the composite risk score and
           the exported report.}
  \label{fig:pipeline}
\end{figure*}


\Cref{fig:pipeline} shows the end-to-end pipeline. After a repository is
registered, \cs\ walks the working tree, applies language filters, and
produces a file inventory. Three analyses then run over that inventory;
the first two are in-process and the third is delegated to the local
model. \Cref{tab:modalities} maps each modality to the weakness
categories it is responsible for.

\begin{table}[t]
  \centering
  \caption{Analysis modalities and the weakness categories each is
           primarily responsible for.}
  \label{tab:modalities}
  \renewcommand{\arraystretch}{1.15}
  \small
  \begin{tabularx}{\columnwidth}{@{}l X@{}}
    \toprule
    \textbf{Modality} & \textbf{Primary responsibility} \\
    \midrule
    Pattern detection &
      Hard-coded secrets, injection sinks, unsafe deserialisation,
      request forgery, private-key material, missing authorisation
      (CWE-77/89/502/918/798)~\cite{mitrecwe}. \\
    Complexity engine &
      Per-function cyclomatic complexity, decision-point counts, and
      structural hot-spot ranking~\cite{mccabe1976}. \\
    LLM semantic review &
      Context-dependent issues: prompt-context handling, unbounded
      resource use, error-suppression, non-determinism, and design
      weaknesses specific to LLM-integrated
      applications~\cite{owaspllm2025}. \\
    Dynamic sandbox &
      Buildability, start-up and resource behaviour, and coarse runtime
      file/socket/process activity~\cite{merkel2014docker}. \\
    \bottomrule
  \end{tabularx}
\end{table}


\subsubsection{Pattern-based detection}
Each source file is scanned with a curated set of syntactic and
lightweight semantic rules targeting high-severity, low-ambiguity
classes: hard-coded secrets and high-entropy literals, command/SQL/prompt
injection sinks, unsafe deserialisation (e.g.\ \texttt{pickle.load} on
non-trusted input), server-side request forgery patterns, private-key
material, and missing-authorisation markers on privileged routes. Rules
carry a severity, a Common Weakness Enumeration
identifier~\cite{mitrecwe} where applicable, and a suggested remediation.

\subsubsection{Cyclomatic complexity engine}
For every function, \cs\ computes McCabe cyclomatic complexity $CC$ from
the control-flow graph~\cite{mccabe1976}, counts decision points, and
records the functions whose $CC$ exceeds a configurable budget as
\emph{structural hot-spots}. The mean complexity $\overline{CC}$ over all
analysed functions feeds the composite score (\cref{sec:risk}).

\subsubsection{On-device LLM semantic review}
\label{sec:llm}
The \textsc{AiService} sends each file (or, for large files, a bounded
window) to the local \mbox{Llama~3.2} model through the Ollama HTTP
endpoint on \texttt{127.0.0.1}. The prompt has three parts: a fixed
system preamble that defines the reviewer role and refusal policy; the
code under review, delimited so that its content cannot be interpreted as
instructions; and a response schema requesting a JSON array of findings
with \texttt{severity}, \texttt{title}, \texttt{line},
\texttt{description}, and \texttt{suggestedFix} fields. Decoding uses a
low temperature and a fixed seed so that repeated audits of an unchanged
file are stable. The returned JSON is validated; malformed items are
discarded; surviving items are tagged as model-originated.

\subsubsection{Dynamic sandbox analysis}
\label{sec:dynamic}
The \textsc{ExecutionService} builds the project into a Docker
image, synthesising a minimal Dockerfile when the repository lacks
one, and runs the container under CPU and memory
limits~\cite{merkel2014docker}. It records image build time, container
start-up latency, and a two-second-interval stream of CPU and memory
utilisation, plus a log of coarse behavioural events (file, socket, and
process activity). Runtime figures are attached to the project and shown
on the dynamic-analysis and container-insight screens.

\subsection{Finding Fusion}
\label{sec:fusion}
Pattern findings and model findings are merged on the tuple
$(\textit{file},\, \textit{normalised line},\, \textit{category})$.
When both modalities report the same location and category, the entry is
kept once and marked \emph{corroborated}; a model finding with no
pattern counterpart is retained and marked \emph{model-only}; a pattern
finding with no model counterpart is retained and marked
\emph{pattern-only}. Severity is taken as the maximum of the two. The
fused list is what every downstream screen and the report consume.

\subsection{Composite Risk Model}
\label{sec:risk}
\begin{figure}[t]
  \centering
  \begin{tikzpicture}[
      font=\footnotesize,
      inp/.style={rounded corners=2.5pt, draw=csslate!45, fill=csgraybg,
                  line width=0.5pt, minimum height=7mm, minimum width=24mm,
                  align=center, font=\scriptsize},
      sub/.style={rounded corners=3pt, draw=csblue!45, fill=csbluebg,
                  line width=0.6pt, minimum height=8mm, minimum width=24mm,
                  align=center, font=\scriptsize},
      fl/.style={-{Stealth[length=1.8mm]}, draw=csslate!60, line width=0.7pt},
    ]

    \node[inp] (nsev) at (0,1.1)  {$n_{\text{sev}}=6$};
    \node[inp] (cc)   at (0,0)    {$\overline{CC}=7.8$};
    \node[inp] (a)    at (0,-1.1) {$a=1$};

    \node[sub] (sv) at (3.2,1.1)  {$s_v = 40$};
    \node[sub] (sc) at (3.2,0)    {$s_c = 61$};
    \node[sub] (sq) at (3.2,-1.1) {$s_q = 95$};

    \node[rounded corners=3pt, draw=csslate!55, fill=white, line width=0.7pt,
          minimum height=12mm, minimum width=20mm, align=center]
          (sum) at (6.4,0) {$H=$\\$0.5s_v$\\$+\,0.3s_c$\\$+\,0.2s_q$};

    \draw[fl] (nsev) -- (sv);  \draw[fl] (cc) -- (sc);  \draw[fl] (a) -- (sq);
    \draw[fl] (sv.east) -- (sum.west);
    \draw[fl] (sc.east) -- (sum.west);
    \draw[fl] (sq.east) -- (sum.west);

    % health-index band
    \begin{scope}[shift={(-0.2,-2.65)}]
      \fill[csredbg,   draw=csslate!30] (0,0)   rectangle (3.4,0.5);
      \fill[csamberbg, draw=csslate!30] (3.4,0) rectangle (6.14,0.5);
      \fill[csgreenbg, draw=csslate!30] (6.14,0) rectangle (6.9,0.5);
      \node[font=\scriptsize, color=csred]   at (1.7,0.25)  {critical $<50$};
      \node[font=\scriptsize, color=csamber] at (4.75,0.25) {moderate 50--79};
      \node[font=\scriptsize, color=csgreen, anchor=west] at (7.0,0.25)
        {low $\ge 80$};
      % case-study marker at H = 57  ->  x = 57/100 * 6.9
      \draw[csink, line width=0.9pt] (3.933,-0.15) -- (3.933,0.65);
      \node[font=\scriptsize\bfseries, color=csink, anchor=south]
        at (3.933,0.65) {$H = 57$};
    \end{scope}
    \draw[fl] (sum.south) |- (2.6,-2.4);

  \end{tikzpicture}
  \caption{Composite risk model. Three disclosed sub-scores
           (\cref{eq:sv,eq:sc,eq:sq}) are combined by a fixed weighting
           into the \phindex\ $H$ (\cref{eq:h}), which is mapped to three
           bands. The marker shows the case-study result, $H=57$.}
  \label{fig:riskmodel}
\end{figure}


\cs\ reports one headline number, the \phindex\ $H \in [0,100]$, computed
from three sub-scores, each also in $[0,100]$ and each disclosed to the
user (\cref{fig:riskmodel}). Let $n_{\text{sev}}$ be the number of fused
findings at \emph{critical} or \emph{high} severity, $\overline{CC}$ the
mean cyclomatic complexity, and let the audit-completion flag be
$a \in \{0,1\}$. Then
\begin{align}
 s_v &= \max\!\left(0,\; 100 - 10\, n_{\text{sev}}\right), \label{eq:sv}\\
 s_c &= \max\!\left(0,\; 100 - 5\, \overline{CC}\right), \label{eq:sc}\\
 s_q &= 95\,a + 50\,(1-a), \label{eq:sq}\\
 H &= \operatorname{round}\!\left(0.5\, s_v + 0.3\, s_c + 0.2\, s_q\right).
 \label{eq:h}
\end{align}
The weights in \cref{eq:h} encode the platform's priorities: exploitable
weaknesses dominate, structural debt is a secondary drag, and analysis
completeness is a small confidence term. $H$ is mapped to three bands:
\emph{low risk} for $H \ge 80$, \emph{moderate risk} for
$50 \le H < 80$, and \emph{critical risk} for $H < 50$. The model is
deliberately simple and transparent; \cref{sec:discussion} discusses
calibration as future work.

\subsection{Reporting}
An audit can be exported as a paginated PDF, executive summary,
structural hot-spot matrix, and a full finding ledger with severity,
type, location, and identifier, or as machine-readable JSON for
downstream tooling. Reports are written to local disk only.

\subsection{Implementation}
\cs\ is implemented in TypeScript on Electron~41 and React~18; charts use
a declarative SVG library and all state is held in an in-process store
backed by SQLite. Optional settings enforce offline mode, disable any
external link handling, and enable encrypted local storage of results.
The prototype is approximately 6\,KLOC of application code excluding
generated UI primitives, and is publicly available~\cite{codesentinelrepo}.
